\chapter{Fazit und Ausblick}
\label{chap:fazit}

Im Rahmen dieser Arbeit konnte gezeigt werden, dass das FM350-2 Modul der SIMATIC S7-300 erfolgreich durch ein VHDL-basiertes Applikationsbeispiel auf dem TM FAST Modul der SIMATIC S7-1500 ersetzt werden kann.  
Alle für Dosieranlagen relevanten Funktionen wurden umgesetzt und in das neue Design integriert.  

Ein zentrales Ergebnis ist die signifikante Verkürzung der Reaktionszeit: Während die FM350-2 Baugruppe eine Verzögerung von rund 50\,µs aufwies, erreicht die neue Implementierung eine Zykluszeit von 20\,ns. Dadurch wird eine sofortige und deterministische Signalverarbeitung ermöglicht, was die Eignung für zeitkritische Dosierprozesse deutlich erhöht.  

Darüber hinaus konnte die Kanalanzahl von ursprünglich 8 auf 14 im TM FAST Modul erweitert werden. Dies eröffnet neue Möglichkeiten für den parallelen Betrieb mehrerer Dosierkanäle und steigert sowohl die Flexibilität als auch die Effizienz der Anwendung.  

Auch die Ressourcennutzung wurde optimiert: Durch gezielte Anpassungen im VHDL-Design konnte die Auslastung der FPGA-Logikelemente auf maximal 80\,\% begrenzt werden, sodass Reserven für zukünftige Erweiterungen bestehen.  

Schließlich wurde das Applikationsbeispiel vollständig in das TIA Portal integriert. Damit steht eine nahtlose Einbindung in bestehende Engineering-Umgebungen zur Verfügung, was die Inbetriebnahme und Diagnose im industriellen Einsatz erheblich erleichtert.  

Insgesamt belegt die Arbeit, dass das TM FAST Modul eine leistungsfähige und zukunftssichere Alternative zum FM350-2 darstellt. Die erfolgreiche Umsetzung der funktionalen und nicht-funktionalen Anforderungen schafft die Grundlage für den Einsatz moderner FPGA-Technologie in hochkanaligen Dosieranwendungen und verdeutlicht deren Potenzial für die industrielle Automatisierung.

\section{Kritische Reflexion und Herausforderungen}

Ein zentrales Problem bestand darin, dass kein realer Aufbau der alten FM 350-2-Baugruppe verfügbar war. Vergleichsdaten über die Zykluszeit mussten daher ausschließlich aus der Dokumentation und dem Handbuch entnommen werden. Dadurch konnten bestimmte Messungen oder Tests zum Vergleich der Reaktionszeiten und Fehlerdiagnosen nicht durchgeführt werden. Diese Einschränkung führte dazu, dass die Analyse an diesen Teilen theoretisch blieb und auf Annahmen basierte, die nicht vollständig verifiziert werden konnten.  

\section{Potenzielle Weiterentwicklungen und zukünftige Anwendungen}

Ein wichtiger Ansatzpunkt für zukünftige Arbeiten ist die Integration von Prozessalarmen, wie sie in der FM 350-2 vorgesehen waren. Diese könnten über eine geeignete Erweiterung des azyklischen Datensatzes \texttt{RD\_REC} realisiert werden, um sowohl detaillierte Fehlerinformationen als auch Diagnosehinweise sicher zu übertragen. Ziel sollte es sein, die Vorteile der neuen Architektur – insbesondere die extrem kurze Reaktionszeit und die sichere Fehlerbehandlung – mit der umfangreichen Diagnosetiefe der alten Baugruppe zu kombinieren.  

Darüber hinaus bietet das flexible VHDL-Design eine solide Grundlage für weitere Applikationen. Durch die Anpassung der Kanalzuordnung und die Erweiterung der Kommunikationsschnittstellen können neue Einsatzfelder erschlossen werden, beispielsweise für hochpräzise Dosier- oder Zählsysteme in unterschiedlichen industriellen Bereichen. Auch die Anbindung an moderne Steuerungssysteme über standardisierte Protokolle stellt ein mögliches Entwicklungsfeld dar, um die Integration in komplexe Automatisierungsumgebungen weiter zu vereinfachen.  
